\chapter{Research Agenda and Open Problems}
\label{ch:frfp-open-problems}

This chapter sketches a research agenda emerging from FRFP and the
case studies in Volume~II.
It is intentionally short and selective:
its goal is to identify a few tractable directions where theory,
empirics, and institutional work can reinforce each other, rather
than to catalogue every possible question.

We group open problems into four clusters:

\begin{enumerate}
  \item foundations and theory;
  \item empirical characterisation of long-horizon systems;
  \item institutional and governance design;
  \item tools and practice for FRFP-aligned development.
\end{enumerate}

\section{Foundations and Theory}

\subsection{T1. Formal conditions for safe horizons}

FRFP treats safe horizons as structural bounds on how far a run in
\(N \ltimes E\) may proceed without re-grounding in tacit judgement.
Volume~I defines this notion; Volume~II shows qualitatively where
horizon failures occur.
We lack sharp theorems of the form:

\begin{quote}
  Given a class of systems with properties \(A,B,C\), any run of
  depth \(>D\) without re-grounding has probability at least \(p\) of
  entering an inadmissible region.
\end{quote}

Open problems:

\begin{itemize}
  \item characterise safe horizons in simplified models:
        bandit processes, Markov decision processes, or repeated
        games with partial monitoring;
  \item prove sufficient conditions under which periodic intervention
        by tacit operators \(T\) guarantees bounds on error or drift;
  \item analyse trade-offs between horizon length, intervention
        frequency, and performance.
\end{itemize}

\subsection{T2. Population-level FRFP dynamics}

Volume~II informally treats populations of users, institutions, and
systems.
A more formal treatment would:

\begin{itemize}
  \item model distributions over tacit states \(T\) (e.g.\ beliefs,
        skills, norms) for populations;
  \item define institutional operators acting on these distributions
        (platform policies, regulations, scientific norms);
  \item study fixed points, cycles, and instability modes when AI
        systems and institutions co-evolve.
\end{itemize}

Concrete questions:

\begin{itemize}
  \item Under what conditions do recommender or LLM-mediated
        environments converge to narrow attractors in belief or
        preference space?
  \item How do different governance operators (e.g.\ content policies,
        audit regimes) alter these dynamics?
\end{itemize}

\subsection{T3. Non-automation theorems at the institutional level}

Volume~I states non-automation results for certain classes of
problems; Volume~II extends this intuition to governance:

\begin{quote}
  There is no purely explicit procedure that can fully replace tacit
  institutional judgement about acceptable socio-technical systems.
\end{quote}

Open problems:

\begin{itemize}
  \item formalise classes of governance tasks (e.g.\ acceptable-risk
        assessment, fairness adjudication) and show limits on their
        automation;
  \item connect these limits to impossibility results in social
        choice, mechanism design, or learning theory;
  \item identify precise boundaries where partial automation is safe
        or beneficial.
\end{itemize}

\subsection{T4. Semantics for large foundation models}

Foundation models and general-purpose LLMs have diffuse explicit
semantics.
We need better theoretical tools to:

\begin{itemize}
  \item describe what these models are actually computing, beyond
        generic sequence prediction;
  \item relate this to task-specific semantics when they are
        fine-tuned or prompted for particular roles;
  \item characterise when and how their behaviour can be constrained
        by FRFP-style boundaries and safe horizons.
\end{itemize}

Open problems include:

\begin{itemize}
  \item compositional semantics for prompted or tool-augmented
        models;
  \item formal models of hallucination and its inevitability under
        certain training regimes;
  \item conditions under which foundation models can safely serve as
        components inside FRFP architectures.
\end{itemize}

\section{Empirical Characterisation of Long-Horizon Systems}

\subsection{E1. Measuring tacit-skill erosion and preservation}

Volume~II highlights tacit erosion as a recurring risk.
Empirical work is needed to:

\begin{itemize}
  \item measure how prolonged AI use affects human skills in law,
        medicine, engineering, and everyday tasks;
  \item disentangle short-term productivity gains from long-term
        dependence;
  \item evaluate design interventions (e.g.\ mandatory verification,
        occasional AI-free tasks) that preserve key skills.
\end{itemize}

This includes controlled experiments, longitudinal field studies, and
instrumentation of real systems.

\subsection{E2. Boundary capture in practice}

Many case studies feature systems that are ``advisory'' on paper but
decisive in practice.
We currently lack good empirical measures of boundary capture.

Open questions:

\begin{itemize}
  \item How often do humans override AI recommendations in different
        domains, and why?
  \item How do override rates evolve over time as trust and habit set
        in?
  \item What interface and incentive designs help maintain meaningful
        boundaries?
\end{itemize}

\subsection{E3. Horizon creep and drift}

Horizon creep---the gradual extension of time, scale, or domain
beyond initial intentions---is visible in several failures, but rarely
measured.

Research directions:

\begin{itemize}
  \item define empirical indicators of horizon creep (expansion of
        user base, jurisdictions, asset classes, patient populations);
  \item correlate these with incident rates, near-misses, and
        qualitative signs of misalignment;
  \item study organisational dynamics that accelerate or resist
        horizon creep (competition, regulation, internal culture).
\end{itemize}

\subsection{E4. Comparative case-study programmes}

Volume~II provides a handful of detailed case studies.
A larger research programme could:

\begin{itemize}
  \item build a curated library of FRFP-analysed systems across
        domains and scales;
  \item develop standardised coding schemes for patterns and
        anti-patterns;
  \item quantitatively test hypotheses such as:
        ``systems with engineered safe horizons have lower incidence
        of catastrophic failure, controlling for domain and scale''.
\end{itemize}

This would turn FRFP from a primarily conceptual framework into a
cumulative empirical research area.

\section{Institutional and Governance Design}

\subsection{G1. FRFP-aware regulatory regimes}

Regulators in finance, healthcare, transport, and digital platforms
increasingly face AI systems with long-horizon effects.
Open questions:

\begin{itemize}
  \item How can regulatory frameworks be rewritten to explicitly
        address:
        \begin{itemize}
          \item task semantics and domain of validity;
          \item boundaries (who decides what);
          \item engineered safe horizons;
          \item institutional non-automation?
        \end{itemize}
  \item What templates for permits, approvals, and ongoing reporting
        best capture these elements?
  \item How can regulators collaborate across domains using a shared
        FRFP vocabulary?
\end{itemize}

\subsection{G2. Organisational FRFP architectures}

Within organisations, there is a design space of governance
architectures:

\begin{itemize}
  \item centralised vs.\ federated AI risk committees;
  \item domain-specific vs.\ cross-cutting boards;
  \item internal vs.\ external oversight mechanisms.
\end{itemize}

Research questions:

\begin{itemize}
  \item Which architectures are most effective at preventing
        anti-patterns (metric worship, advisory-in-name-only, horizon
        creep)?
  \item How can FRFP concepts be embedded in standard corporate
        processes (risk registers, audits, product reviews)?
  \item How do different incentive structures affect the success of
        FRFP-style governance?
\end{itemize}

\subsection{G3. Professional norms and FRFP}

Many domains (medicine, law, science, engineering) already have strong
professional norms around judgement, evidence, and responsibility.

Open problems:

\begin{itemize}
  \item How can FRFP concepts be translated into domain-native
        language and integrated into existing codes of practice?
  \item What training and certification regimes are needed for
        ``FRFP-literate'' professionals and leaders?
  \item How can professional societies act as FRFP operators,
        updating norms as AI capabilities change?
\end{itemize}

\section{Tools and Practice}

\subsection{P1. FRFP-aware development and operations tools}

Today’s MLOps and product tools are largely agnostic to FRFP
structure.

Opportunities:

\begin{itemize}
  \item development tools that:
        \begin{itemize}
          \item encourage explicit task-definition and scope;
          \item annotate model outputs with uncertainty and provenance;
          \item make boundary events first-class objects.
        \end{itemize}
  \item operations tools that:
        \begin{itemize}
          \item track safe horizons and trigger reviews;
          \item surface FRFP metrics to governance bodies;
          \item log institutional decisions as part of the system
                state.
        \end{itemize}
\end{itemize}

\subsection{P2. Simulation environments for FRFP systems}

Several case studies depend on long-horizon, population-level
effects.
We need simulation tools that:

\begin{itemize}
  \item model interactions between AI systems, users, and institutions;
  \item allow testing of governance and boundary designs before
        deployment;
  \item support ``what if'' analysis:
        what happens if we change the ODD, screening interval, or
        content policy?
\end{itemize}

This suggests a line of work on FRFP-aware agent-based simulations and
testbeds.

\subsection{P3. Education, documentation, and patterns}

Finally, there is practical work to do in:

\begin{itemize}
  \item creating FRFP teaching materials for engineers, policy-makers,
        and students;
  \item collecting and publishing patterns and anti-patterns (as in
        Chapter~\ref{ch:frfp-patterns}) in a form usable by
        practitioners;
  \item developing documentation standards that include FRFP elements
        (semantics, boundaries, horizons, governance) alongside
        conventional API or model docs.
\end{itemize}

\section{Closing Remarks}

FRFP was motivated by a gap:
the absence of a formal language for reasoning about long-horizon
human--AI systems that are simultaneously technical, social, and
institutional.
Volume~I filled part of that gap at the theoretical level;
Volume~II tested the language against concrete cases.

The research agenda outlined here points in both directions:

\begin{itemize}
  \item inward, toward deeper formal results on safe horizons,
        population dynamics, and non-automation;
  \item outward, toward empirical measurement, institutional design,
        and practical tools.
\end{itemize}

The hope is that FRFP can function as a common scaffold across these
efforts: a way for theorists, engineers, clinicians, financiers,
regulators, and others to talk about the same objects---explicit
spaces, tacit spaces, boundaries, horizons, and operators---even as
they work at very different levels of abstraction and practice.
